% =====================================================================
\section{The exact spectral squeeze, and why H1 is necessary but not sufficient}

\figladder
\label{sec:squeeze}
% =====================================================================

\subsection{Setup and notation}

Let $\nu\ge0$ be the K\"all\'en--Lehmann measure of $\operatorname{tr}F^2$, with
$\operatorname{supp}\nu\subset[m_1^2,\infty)$, atom weight $f_1^2=\nu(\{m_1^2\})$, and moments
$M_k=\int s^k\,d\nu(s)$. In practice $M_0=\int d\nu$ is the total spectral weight of the
operator in the channel, and $M_1$ is the first moment (the ``average mass squared'' weighted by
residue). Let $\xi$ be the zero-momentum coarse mode
($\xi=\langle v,\log U\rangle$ for the relevant mode; $\operatorname{Var}(\xi)$ is the
zero-momentum connected propagator). The \textbf{susceptibility} is
\begin{equation}
\chi \;:=\; \operatorname{Var}(\xi) \;=\; \int s^{-1}\,d\nu(s)
\qquad(\text{the zero-momentum connected two-point function}).
\label{eq:chi}
\end{equation}
This is the inverse of the effective curvature only in the Gaussian case; in general it is a
residue-weighted \emph{harmonic} mean of the masses, dominated by the lightest residue. It is
not the lightest mass. The equality of the two expressions in \eqref{eq:chi} is the
zero-momentum Ward statement, up to the normalization of $O$; the squeeze below uses only ratios
and is normalization-free.

We record three elementary facts about $\chi$ that will be used repeatedly, and that explain why
it is the ``wrong'' functional for a lower bound on the mass.

\begin{enumerate}
\item \textbf{Scale.} $\chi$ has the dimension of an inverse mass squared. It is large when
  there is light spectral weight (a small $s$ with non-negligible residue) and small when the
  weight sits high.
\item \textbf{Harmonic mean.} If $\nu$ is a probability measure on masses, then
  $\chi=\mathbb{E}_\nu[s^{-1}]$, the harmonic mean of $s$. The harmonic mean is dominated by the
  \emph{smallest} values, but it is dominated by their \emph{residues}, not by their location:
  a tiny residue at tiny mass contributes $w/\varepsilon$, which can be made arbitrarily small by
  taking $w\to0$ at fixed $\varepsilon$, or kept fixed by taking $w\sim\varepsilon$.
\item \textbf{No lower edge.} $\chi$ is continuous under weak-$*$ convergence of measures against
  bounded continuous test functions on $(0,\infty)$; it does not see the location of the bottom
  of the support except through its residue-weighted contribution. The bottom atom can be pushed
  to zero while $\chi$ stays bounded (the thin-atom proposition below).
\end{enumerate}

\subsection{The squeeze \LT}
\label{sec:squeeze-thm}

For any such $\nu$:
\begin{equation}
\frac{f_1^2}{\chi}\ \le\ \inf\operatorname{supp}\nu\ \le\ \frac{M_0}{\chi}.
\label{eq:squeeze}
\end{equation}
\begin{proofx}
\textbf{Lower inequality.} $\chi m_1^2=f_1^2+m_1^2\int_{s>m_1^2}s^{-1}d\nu(s)\ge f_1^2$, since
$m_1^2s^{-1}\ge1$ on the support above $m_1^2$. Hence $m_1^2\ge f_1^2/\chi$.

\textbf{Upper inequality.} On $\operatorname{supp}\nu$, $s\ge m_1^2$, so $s^{-1}\le m_1^{-2}$,
hence $\chi=\int s^{-1}d\nu\le M_0/m_1^2$, i.e.\ $m_1^2\le M_0/\chi$.
\end{proofx}

Both inequalities are elementary and exact. They are the entire content of the ``squeeze'' and
they contain no dynamical input. Three remarks are worth making precise.

\textbf{(a) Equality cases.} The lower inequality is an equality iff $\nu$ has no mass above
$m_1^2$, i.e.\ $\nu=f_1^2\delta_{m_1^2}$ (a single atom; then $\chi=f_1^2/m_1^2$). The upper
inequality is an equality iff all mass sits at $m_1^2$ (the same degenerate case). For any
non-degenerate measure the two inequalities are strict; in particular, for a measure with an atom
at $m_1^2$ and a continuum above, both are strict and the true bottom lies strictly between the
two bounds.

\textbf{(b) The bounds are sharp over all measures.} Fix $M_0=1$. For every $\eta>0$ there is a
measure with $f_1^2/\chi\le\eta$ (thin atom at tiny mass, the thin-atom proposition below) and a
measure with $M_0/\chi\ge1/\eta$ (mass concentrated just above zero with small residue, so
$\chi$ is large). Hence no improvement of either constant is possible without extra structure.

\textbf{(c) What the squeeze does \emph{not} say.} It does not say that $\chi$ bounds the mass
from below, and it does not say that a large $\chi$ implies a light state. It says that the mass
is squeezed between the residue-to-susceptibility ratio and the total-weight-to-susceptibility
ratio. Both ratios involve $\chi$ in the denominator, so an \emph{upper} bound on $\chi$ is an
\emph{lower} bound on $m_1^2$ only if the numerator is also controlled from below. This is the
entire reason the campaign's missing input is a pair, not a single inequality.

\subsection{H1 is necessary, not sufficient}
\label{sec:h1-notsuff}

The condition H1 is
\begin{equation}
\mathrm{H1}:\qquad \operatorname{Var}_\mu(\xi_{k\to0})=\chi\ \le\ 1/\lambda,
\qquad\text{uniform in }(L,a).
\label{eq:h1}
\end{equation}

\textbf{Necessity \LT.} If the theory has a gap $m_1\ge\sqrt\lambda$ (with the normalization
$M_0$ fixed), then $\chi=\int s^{-1}d\nu\le M_0/\lambda$, so $\chi\le M_0/\lambda$. This is a
theorem: $\text{gap}\Rightarrow\chi\le M_0/\lambda$ (up to the normalization constant). The
corpus records it as $\chi=\int s^{-1}d\nu\le M_0/\lambda$ (\code{REDUCTION-AUDIT.md}).

\textbf{Non-sufficiency \LT.} Take (Figure~\ref{fig:measure}, right panel)
\begin{equation}
\nu=w\,\delta_\varepsilon+(1-w)\,\delta_{s_2},\qquad w/\varepsilon\le1/\lambda .
\label{eq:thinatom}
\end{equation}
Then $\chi=w/\varepsilon+(1-w)/s_2<\infty$, so H1 holds, while
$\inf\operatorname{supp}\nu=\varepsilon\to0$. \textbf{H1 holds and the gap fails.} The measure has
a thin atom at $\varepsilon$, invisible to the susceptibility because the susceptibility is a
harmonic mean: a tiny residue at tiny mass contributes $w/\varepsilon$, which can be kept below
$1/\lambda$ by taking $w$ small. The susceptibility sees the \emph{residue-weighted} mass, not the
lightest mass.

\begin{claimbox}{cT}
\claimhead{Proposition}{(thin atom)} \emph{Fix any $\lambda>0$ and any $M_0>0$. There is a
probability measure $\nu$ on $[0,\infty)$ with}
\[
M_0(\nu)=M_0,\qquad \chi(\nu)=\int s^{-1}d\nu\le\frac1\lambda,\qquad
\inf\operatorname{supp}\nu=\varepsilon
\]
\emph{for every $\varepsilon\in(0,\varepsilon_0)$ with $\varepsilon_0=\varepsilon_0(\lambda,
M_0)>0$. In particular, for every $\eta>0$ there is a measure satisfying H1 with
$\inf\operatorname{supp}\nu<\eta$ and $f_1^2=\nu(\{\inf\operatorname{supp}\nu\})=w$ arbitrarily
small.}
\begin{proofx}
Normalize $M_0=1$ (scale $s\mapsto M_0 s$, which rescales $\chi$ by $1/M_0$). Take
$\nu_\varepsilon=w\delta_\varepsilon+(1-w)\delta_{s_2}$ with $w=\varepsilon/(2\lambda)$ and
$s_2=2\lambda(1-w)$. Then $\chi=w/\varepsilon+(1-w)/s_2=1/(2\lambda)+1/(2\lambda)=1/\lambda$,
$M_0=1$, and $\inf\operatorname{supp}\nu=\varepsilon$. The atom weight is $w=\varepsilon/(2\lambda)\to0$.
\end{proofx}
\end{claimbox}
\label{prop:thinatom}

This counterexample is not a pathology; it is the exact obstruction the campaign met
everywhere. It is why the audit's verdict is that the chain written as $\iff$ is unsound: read
against the corpus's own definitions, it is \textbf{(iv) false $\to$ (iii) conjecture $\to$
(ii) definition}, plus one genuine one-directional theorem. The final $\iff$ (H1 $\iff$ gap)
holds in one direction only.

\figmeasure

\textbf{The structural reading.} H1 is an \emph{upper} condition on a second-moment functional.
The atom is the \emph{bottom} of the support. An upper bound on a moment can never force a
bottom: this is the ``upper-bound machine, lower-bound target'' mismatch that recurs in every
wave (Chapter~\ref{ch:immune}, family 14). The thin-atom measure is the exact witness, and it
satisfies every positivity and moment condition that a free-field-blind argument could use.

\subsection{Machine-checked: H1 does not imply an atom}
\label{sec:leanblind}

The logical content of \S\ref{sec:h1-notsuff} is now machine-checked, independent of any
analytic model. The module \code{lean/YMStatement/H1Sharpening.lean} (Lean~4, toolchain pinned;
zero \code{sorry}; axioms only \code{[propext, Classical.choice, Quot.sound]}) proves the
following three statements about real numbers.

\begin{claimbox}{cT}
\claimhead{Theorem}{(Lean \code{lemmaG\_twoSided})} \emph{If $0<m$, $0<\chi$, $f\le\chi m$ and
$\chi m\le M$, then $f/m\le\chi$ and $\chi\le M/m$.}
\end{claimbox}

\begin{claimbox}{cT}
\claimhead{Theorem}{(Lean \code{cert\_necessity})} \emph{If $0<f$, $0<\chi$, $0<\lambda$,
$\chi\le1/\lambda$, $f\le\chi m$ and $0<m$, then $\lambda f\le m$.}
\end{claimbox}

\begin{claimbox}{cT}
\claimhead{Theorem}{(Lean \code{H1\_does\_not\_imply\_atom})} \emph{There is no way to derive
$0<f$ from the hypotheses}
\[
0<\lambda,\quad 0<\chi,\quad \chi\le1/\lambda,\quad 0<m,\quad f\le\chi m,\quad \chi m\le M .
\]
\emph{The tuple $(f,m,M,\chi,\lambda)=(0,1,1,1,1)$ is a countermodel
(\code{H1\_blind\_to\_thin\_atom}).}
\end{claimbox}
\label{thm:leanblind}

\textbf{What this certifies, and what it does not.} The theorems are order-algebraic: they
encode the squeeze as the hypotheses $f\le\chi m$, $\chi m\le M$ and the ceiling as
$\chi\le1/\lambda$. Together they say: \emph{the ceiling plus the squeeze gives the conditional
floor $m\ge\lambda f$, and the only missing ingredient is $f>0$.} In other words,
\[
\mathrm{H1}\ \wedge\ (f_1^2>0)\ \Longrightarrow\ m_1^2\ \ge\ \lambda f_1^2\ >\ 0,
\]
and the second conjunct cannot be manufactured from the first. The formal statement is about the
\emph{logic} of the reduction, not about Yang--Mills; the analytic input $f\le\chi m$ is itself
the squeeze, which is a theorem for the K\"all\'en--Lehmann measure (and is supplied as a
hypothesis in the formalization). The value of the machine check is that the ``H1 is blind to a
thin atom'' claim is no longer an informal remark: it is a compiled countermodel. This is the
first place where the campaign's honesty line is machine-load-bearing.

\textbf{The corresponding sufficiency check.} The module also proves \code{cert\_sufficiency}:
if $0<M$, $0<\chi$, $0<\lambda$, $\chi m\le M$ and $\lambda M\le m$, then $\chi\le1/\lambda$.
That is, a gap floor at the level of $M$ implies the certificate. Together the two directions say
that the certificate is \emph{equivalent} to a gap floor only when the atom weight is fixed and
positive; without the atom it is strictly weaker. The Lean statements are reproduced in full in
Appendix~\ref{app:lean}.

\subsection{The correct sufficient condition}

The repair is the pair:
\begin{equation}
(f_1^2\ge r_0)\ \wedge\ (\chi\le1/\lambda)
\qquad\Longrightarrow\qquad
m_1^2\ \ge\ f_1^2/\chi\ \ge\ r_0\lambda\ >\ 0 .
\label{eq:sufficient}
\end{equation}
This is immediate from the squeeze and it is the campaign's central logical content: \textbf{a
lower bound on the mass needs both a residue floor and a susceptibility ceiling.} The residue
floor is the hard half; the susceptibility ceiling is H1. The campaign's honest statement of the
missing input is therefore not H1 alone but
\begin{equation}
H^*:\qquad f_1^2\ge r_0>0\ \ \text{with isolation } I_2,
\qquad(\text{and }\chi\le1/\lambda\text{ as the matching ceiling}),
\label{eq:hstar}
\end{equation}
where $I_2$ is $\Lambda_{\star}:=\inf\operatorname{supp}(\nu-f_1^2\delta_{m_1^2})$,
$\Lambda_{\star}-m_1^2\ge\delta>0$. The residue--level-gap lemma (Chapter~\ref{sec:theorems})
turns I1 into I2 and back under a self-certifying condition, so $H^*=\text{H1}$ \emph{as a
structural identification}; but the identification is the conjecture, and the residue datum is
strictly stronger than the bare gap.

\textbf{Cross-qualification on $I_2$.} If $I_2$ is consumed through the window theorem
$\rho^2\le C^2$ (Chapter~\ref{sec:theorems}), its constant must be $C=2$, not $1.5$, and the floor it
yields is $m_1^2\ge\Delta/3$ with $\Delta:=m_2^2-m_1^2$ (not $\Delta/1.25$). Since
$\rho^2\le C^2\iff m_1^2\ge t/(1+C^2)$ with $t=T_2(V)$ a computable rank-2 trace, $I_2$ \emph{inherits}
H1 and cannot be used independently of H1. \LN

\textbf{Optimality of the pair.} Neither half alone suffices, and the pair is not improvable in
the following sense. If $f_1^2\ge r_0$ but $\chi$ is unbounded, then $m_1^2$ can be arbitrarily
small: the squeeze only gives $m_1^2\ge f_1^2/\chi$, and $\chi\to\infty$ kills the bound. If
$\chi\le1/\lambda$ but $f_1^2$ is not controlled, the thin-atom construction gives
$m_1^2\to0$. So the pair is exactly the minimal sufficient datum of ``residue times inverse
susceptibility'' type. Any stronger result must use more than the two moments, which is precisely
what the weak-$*$ coordinate (Chapter~\ref{ch:weakstar}) says is impossible for finitely
computable functionals.

\subsection{Why H1 is blind: the convex-order obstruction}
\label{sec:schur}

H1 is an \textbf{upper-only} condition. Every functional of $\nu$ built from its moments or from
positivity is a Schur-concave functional of the measure, hence gives an \textbf{upper} bound on
the atom weight, never a lower bound. We make this precise in two steps.

\textbf{Step 1: the free field is a feasible atomless point.} The massless free field is a
legitimate OS theory with $\nu$ supported on $[0,\infty)$ and $f_1^2=\nu(\{0\})=0$. It satisfies
H1 for every $\lambda>0$ (its $\chi$ is the free inverse propagator). Hence the feasibility set of
measures satisfying H1 (with fixed $M_0$) contains an atomless element. Any inequality of the form
$f_1^2\ge\Phi(\text{moments})$ with $\Phi$ a function of finitely many moments must hold at that
element, forcing $\Phi\le0$ on the feasible set. \textbf{Every such bound is one-sided.}

\textbf{Step 2: finite data cannot distinguish.} Let $K\in\mathbb{N}$ and let $\nu$ be any
measure with finite moments up to $K$. Two constructions are relevant.
\begin{itemize}
\item \emph{Exact-moment atomless replacement.} The first $K+1$ moments of $\nu$ determine a
  positive measure with at most $K+1$ support points (Gaussian quadrature); that measure has
  atoms, but they sit at points determined by the moments and need not include the bottom of
  $\operatorname{supp}\nu$.
\item \emph{Smoothed replacement.} Replacing each quadrature atom $\delta_{x_j}$ by the uniform
  measure on $[x_j-h,x_j+h]$ keeps positivity, keeps the total mass, and moves every moment
  $M_k$, $k\le K$, by $O(h)$. The smoothed measure is atomless, hence has $f_1^2=0$ at its
  bottom, while matching the original data to $O(h)$.
\end{itemize}
Taking $h\to0$ gives: \textbf{for every finite list of moments and every $\eta>0$ there is an
atomless measure whose moments differ from the original by at most $\eta$.} This is the precise
sense in which ``finite data is blind to the atom''; it is the measure-theoretic core of the
weak-$*$ statement of Chapter~\ref{ch:weakstar}, and it is why no moment, positivity, or
convex-order argument can lower-bound $f_1^2$.

\textbf{Convex order.} The phrase ``Schur-concave'' can be made precise. For probability
measures on $\mathbb R$ with a common mean, $\nu\preceq_{\mathrm{cx}}\mu$ means
$\int\varphi\,d\nu\le\int\varphi\,d\mu$ for every convex $\varphi$ (equivalently, equal means
plus the stop-loss inequalities). The Dirac mass $\delta_m$ is the \emph{smallest} element in
convex order at mean $m$, since $\varphi(m)\le\int\varphi\,d\nu$ for convex $\varphi$; every
mean-preserving spread is larger. The variance is a convex-order functional. So an upper bound on
the variance --- which is what H1 is --- is a convex-order condition, and convex order is blind to
where the mass sits: a two-point measure with a thin atom at the bottom and a heavy atom above
can have the same mean and a smaller variance than a concentrated measure at the mean. No
convex-order inequality can force a bottom atom. This is the theorem-level content of
``Schur-concave'', and it is why the search for a forcing tool had to leave the convex-order
class entirely.

\textbf{Consequence for the campaign.} The master obstruction is not a failure of imagination; it
is a structural theorem about the class of arguments. Positivity, reflection positivity, spectral
positivity, analyticity, and every convex-order inequality give \textbf{upper bounds and the
vacuum}, never the edge. This is why the paper repeatedly says ``positivity is blind to the
edge,'' and it is why the search for a forcing tool had to leave the class of bounded-kernel
integral functionals entirely (Chapter~\ref{ch:tools}, the Hardy escape).

\subsection{$G''$ versus $1/\chi$}

The curvature $G''=V_{\mathrm{eff}}''$ and the inverse susceptibility $1/\chi$ coincide only for
a Gaussian marginal. The exact one-directional statement is:
\begin{equation}
\text{the marginal of } \xi \text{ is Gaussian}
\qquad\Longrightarrow\qquad
G''\cdot\operatorname{Var}(\xi)=1 ,
\label{eq:ward}
\end{equation}
where $\xi$ is the coarse mode, so that $G''=\operatorname{Var}(\xi)^{-1}$ in the Gaussian case.
The converse is not available as an identity: for a
non-Gaussian marginal the product can be above or below one. Ward's identity $E[\xi\,f]=1$ (with
$f=\langle v,\nabla S\rangle$) is exact for every measure and every $\beta$; it says that $\xi$ and
$f$ have unit correlation, not that the conditional variance is $G''$. Measured at $\beta=2.5$,
the two differ by $2$--$5\%$ ($G''\approx0.281$ versus $1/\operatorname{Var}\approx0.272$;
\code{FORGE5-SL-IR-SINGLE.md}, quoted in \code{docs/new-tools-registry.md}). So even the
``curvature face'' is not literally the susceptibility; it is a second-moment object that equals
the susceptibility only in the Gaussian limit. The unification $G''=1/\chi=\text{H1}$ is therefore
a \emph{conjecture with a measured $2$--$5\%$ gap}, not an identity. The exact identity that does
hold is the variance decomposition of Chapter~\ref{sec:theorems}.

\subsection{The verdict on the reduction}

The honest logical skeleton (\code{REDUCTION-AUDIT.md} \S``What actually survives'',
Figure~\ref{fig:diamond}) is:

\begin{itemize}
\item \LT{} isolated atom at $\inf\operatorname{supp}\nu$ of a channel $\Rightarrow$ exponential
  clustering in that channel;
\item \LT{} $\text{gap}\Rightarrow\chi\le M_0/\lambda$;
\item \LT{} $G''\ge\lambda\Rightarrow\operatorname{Var}(\xi)\le1/\lambda$ (Brascamp--Lieb, one way);
\item \LC{} the unification $H^*=\text{H1}$ (``one object'') --- a structural identification,
  not a theorem, and now \emph{doubted}: H1 is the inverse susceptibility $1/\chi$ (an upper bound
  on $m_1^2$, from $f_1^2/\chi\le m_1^2\le M_0/\chi$), equal to the pole only for elementary
  single-atom fields; the free massive composite $\mathord{:}\varphi^2\mathord{:}$ has
  $G''=1/\chi\ne m_1^2$;
\item \textbf{(iv) false} $\text{gap}\iff\text{atom}$, and $\text{H1}\iff\text{gap}$.
\end{itemize}

Written as $\iff$, the chain is unsound. Its true shape is a \textbf{diamond of implications
converging on the necessary condition $\chi\le1/\lambda$}, with the sufficient condition being
the strictly stronger $H^*=\text{I1}\wedge\text{I2}$. Proving H1 alone does not yield the gap
(the tiny-residue counterexample, now machine-checked); proving $H^*$ does, but $H^*$ is not H1.
This is the first place where the campaign's honesty line is load-bearing, and we state it
plainly: \textbf{the reduction relocates the residue content; it does not remove it.}
